\chapter{Local and global grains}
\label{chap:grains}

Throughout the structural argument we assume that a positive critical
exponent $\sigma<1$ exists. The WZ1 statements in this chapter are
recorded with their original normalization. Their WZ2 counterparts
use Proposition~\ref{prop:sticky-structure}; the principal output is
the following proposition.

\begin{proposition}[Lipschitz grains; WZ2, Proposition 6.3]
\label{prop:lipschitz-grains}
\lean{Kakeya.Assouad.PureWZ2GrainsFromCriticalStatement,Kakeya.Assouad.pure_wz2_node04_grains}
\leanok
\uses{def:extremal,def:ad-set,def:plane-map}
For every $\eps>0$ and $\delta_0>0$, there is an $\eps$-extremal
configuration at some $0<\delta\leq\delta_0$ with both of the
following properties on the same shaded union $E$:
\begin{enumerate}
\item There is a $1$-Lipschitz function $f:[-1,1]\to\R$ such that
\[
       (E\cap\{z=z_0\})\cdot(1,f(z_0),0)
\]
is a $(\delta,1-\sigma,\delta^{-\eps})_1$-ADset for every
$z_0\in[-1,1]$.
\item There is a $1$-Lipschitz plane map $V:E\to S^2$ such that,
for $p\in E$ and $\delta\leq\rho\leq1$,
\[
       V(p)\cdot(E\cap B(p,\sqrt\rho))
\]
is a $(\rho,1-\sigma,\delta^{-\eps})_1$-ADset.
\end{enumerate}
\end{proposition}
\begin{proof}
\uses{prop:sticky-structure,lem:weak-plane-map,lem:lipschitz-plane-map,
  lem:local-grains-all-scales,lem:mild-rescaling}
\leanok
The multilinear estimate first supplies a weak plane map. Passing to
coarse tubes improves the angular tolerance to the tube radius. The
comparison of plane maps at two scales gives a modulus of continuity;
a finite iteration produces a Lipschitz plane map. Apply the local
covering estimates at successive scales and keep one common
refinement. Finally, choose a coordinate patch where one horizontal
component of the normal is bounded away from zero. The ratio of the
other horizontal component to it supplies the slope $f$. A mild
rescaling absorbs the remaining Lipschitz constants and gives the
global horizontal-slice estimate together with the local one.
\end{proof}

\begin{proposition}[The WZ1 grain proposition; WZ1, Proposition 4.1]
\label{prop:wz1-grains}
\lean{Kakeya.Assouad.WZ1PlaninessGraininessConclusion}
\uses{def:wz1-extremal,def:plane-map,def:ad-set}
At arbitrarily small scales there are WZ1 $\eps$-extremal families
with the two grain properties in Proposition~\ref{prop:lipschitz-grains}:
global horizontal grains with a $1$-Lipschitz slope, and local grains
with a $1$-Lipschitz plane map. The critical exponent is now the WZ1
exponent $\sigma_3$.
\end{proposition}
\begin{proof}
\uses{lem:weak-plane-map,lem:lipschitz-plane-map,lem:local-grains-all-scales}
These are the two conclusions assembled at the end of WZ1 Section 4.
The argument is also the source of the WZ2 proposition above.
\end{proof}

\begin{remark}[The domain of the plane map]
\label{rem:plane-domain}
WZ1 Definition 4.1 defines $V$ on $E$, but the local-grain display in
WZ1 Proposition 4.1 and WZ2 Propositions 6.3--6.5 quantifies over
$p\in\R^3$. Without specifying an extension, $V(p)$ is undefined
outside $E$. The corresponding Lean object,
\texttt{PureWZ2LocalGrainData}, quantifies over $p\in E$. When a later
argument uses a geometric cell centre, it first selects an occupied
point in that cell and compares the two neighbourhoods. The statement
above uses this precise domain.
\end{remark}

\section{From weak planiness to a continuous plane field}

\begin{lemma}[Weak planiness; WZ1, Lemma 4.3]
\label{lem:weak-plane-map}
\lean{Kakeya.Assouad.WZ1WeakPlaninessAssemblyConclusion}
\uses{def:extremal}
For sufficiently small source loss $\eta$ and fine scale $\delta$,
an $\eta$-extremal configuration has a refinement and a unit normal
field $V$ satisfying
\[
       |V(p)\cdot\dir(T)|\leq\delta^{\sigma/2}
       \qquad(p\in Y'(T)).
\]
\end{lemma}
\begin{proof}
\uses{thm:multilinear-kakeya,lem:robust-transversality}
The multilinear estimate bounds the measure of points supporting many
triples with large determinant. Such points carry little shaded mass
in an extremal configuration. At each remaining point choose two
directions separated by a definite power of $\delta$, using robust
transversality. Their cross product defines a normal. The determinant
bound puts all retained directions close to its perpendicular plane.
\end{proof}

\begin{lemma}[A coarse plane map; WZ1, Lemma 4.4]
\label{lem:coarse-plane-map}
\lean{Kakeya.Assouad.WZ1CoarsePlaneMapConclusion}
\uses{def:extremal,def:plane-map}
Given $\eps>0$, choose $\eta$ and $\delta$ sufficiently small. Suppose
an $\eta$-extremal family has a unit normal field $V$ with
$|V(p)\cdot\dir(T)|\leq\rho/2$ on its shaded incidences, where
$\delta^{1-\eps}\leq\rho\leq\delta^\eps$.
There is a refined fine family covered by an $\eps$-extremal
$\rho$-tube configuration with a plane map $\widetilde V$ constant
on each shaded $\rho$-cube. Its value on each such cube is the value
of $V$ at a retained point of that cube.
\end{lemma}
\begin{proof}
\uses{prop:sticky-structure}
Take a balanced cover and, in each coarse cube, choose a fine point
with maximal surviving multiplicity. Retain the coarse incidences
witnessed by its fine tubes. Directional closeness transfers the
$\rho/2$ bound to a radius-$\rho$ plane-map bound. The balance and
multiplicity estimates ensure that the surviving coarse shading is
still dense enough to be extremal.
\end{proof}

\begin{lemma}[Oscillation in a cube; WZ1, Lemma 4.5]
\label{lem:plane-oscillation}
\lean{Kakeya.Assouad.WZ1TwoScalePlaneMapConclusion}
\uses{def:plane-map,def:extremal}
For every $\eps>0$, an $\eta$-extremal family with a plane map $V$
and an intermediate scale
$\rho\in[\delta^{1-\eps},\delta^\eps]$ has an $\eps$-extremal
refinement on which
\[
        |V(p)-V(q)|\leq\rho
\]
whenever $p,q$ belong to one $\rho$-cube. Here $\eta$ and then
$\delta$ are chosen sufficiently small in terms of $\eps$.
\end{lemma}
\begin{proof}
\uses{lem:coarse-plane-map,lem:robust-transversality}
Compare the normals through two nested balanced covers. A pair of
well-separated directions determines a plane to within the required
error. The close-direction count excludes the pairs for which this
determination is unstable. Refining to the surviving incidences gives
the common-cube estimate.
\end{proof}

\begin{corollary}[Oscillation in a ball; WZ1, Corollary 4.6]
\label{cor:plane-oscillation}
\lean{Kakeya.Assouad.WZ1ExtremalOneScalePlaneMapConclusion}
\uses{def:extremal,def:plane-map}
Under the same hypotheses, there is an $\eps$-extremal refinement
such that $|V(p)-V(q)|\leq2\rho$ whenever $|p-q|\leq\rho$.
\end{corollary}
\begin{proof}
\uses{lem:plane-oscillation}
Apply the cube estimate to a bounded collection of shifted grids and
retain a common refinement. Nearby occupied points can then be
compared in one of the enlarged cubes.
\end{proof}

\begin{lemma}[Finite-scale Lipschitz refinement; WZ1, Lemma 4.7]
\label{lem:finite-lipschitz}
\lean{Kakeya.Assouad.WZ1FinitePlaneMapIterationConclusion}
\uses{def:extremal,def:plane-map}
For every $\eps>0$, a sufficiently extremal family whose plane map
is constant on $\delta$-cubes has an $\eps$-extremal refinement on
which that plane map is $\delta^{-\eps}$-Lipschitz.
\end{lemma}
\begin{proof}
\uses{cor:plane-oscillation}
Apply the preceding corollary at a fixed finite sequence of power
scales. Choose the losses in reverse order, so that each output meets
the hypotheses of the next application. For two points use the first
scale above their separation. The gap between consecutive scales
accounts for the final factor $\delta^{-\eps}$; constancy on the
finest cells handles the bottom scale.
\end{proof}

\begin{lemma}[Existence of a Lipschitz plane map; WZ1, Lemma 4.8]
\label{lem:lipschitz-plane-map}
\lean{Kakeya.Assouad.WZ1LipschitzPlaneMapConclusion}
\uses{def:plane-map,def:extremal}
For every $\eps,\delta_0>0$, there is an $\eps$-extremal family at
some $0<\delta\leq\delta_0$ with a $\delta^{-\eps}$-Lipschitz plane
map.
\end{lemma}
\begin{proof}
\uses{lem:weak-plane-map,lem:coarse-plane-map,lem:finite-lipschitz}
Obtain a weak plane map, pass to a scale at which it is a genuine
plane map constant on the coarse cubes, and apply the finite-scale
refinement. Renaming the resulting scale gives the statement.
\end{proof}

\section{The local covering estimates}

In the next three statements the source pair is $\eta$-extremal and
has a $\delta^{-\eta}$-Lipschitz plane map. For each prescribed
$\eps>0$, $\eta$ and the fine-scale threshold are chosen afterwards.
The notation $V(Q)$ means a value at an occupied point of $Q$; the
Lipschitz bound controls the effect of the choice.

\begin{lemma}[Projection of one cube; WZ1, Lemma 4.9]
\label{lem:grain-cube}
\lean{Kakeya.Assouad.WZ1LocalGrainCubeCountConclusion}
\uses{def:plane-map,def:extremal}
For $\rho\in[\delta^{1-\eps},\delta^\eps]$ and
$\tau\in[\rho^{1-\eps},\sqrt\rho]$, there is an $\eps$-extremal
refinement such that, for every occupied $\tau$-cube $Q$,
\[
 \cE_\rho\bigl(V(Q)\cdot(E_{Y'}\cap Q)\bigr)
       \leq\rho^{-\eps}(\tau/\rho)^{1-\sigma}.
\]
\end{lemma}
\begin{proof}
\uses{prop:sticky-structure,lem:robust-transversality}
Use a balanced cover at scale $\rho$ and count the coarse tubes that
meet $Q$. Local planiness identifies their normal projections up to
scale $\rho$. A large number of distinct projected intervals would
force too much shaded volume, using the lower bound on each
rescaled fibre. The extremal volume upper bound gives the estimate.
\end{proof}

\begin{lemma}[Localization to an interval; WZ1, Lemma 4.10]
\label{lem:grain-interval}
\lean{Kakeya.Assouad.WZ1LocalGrainIntervalConclusion}
\uses{def:extremal,def:plane-map}
At the same scales, an $\eps$-extremal refinement can be chosen so
that, if $Q$ is an occupied $\sqrt\rho$-cube and
$F=V(Q)\cdot(E_{Y'}\cap Q)$, then every interval $J$ of length
$\tau$ satisfies
\[
                       \cE_\rho(F\cap J)
                       \leq\rho^{-\eps}(\tau/\rho)^{1-\sigma}.
\]
\end{lemma}
\begin{proof}
\uses{lem:grain-cube}
Localize the shaded incidences above $J$, and cover the relevant
portion of $Q$ by cubes at the smaller scale. The plane-map variation
changes the projected interval by at most a bounded enlargement.
Apply the one-cube estimate and retain a common refinement over the
required intervals.
\end{proof}

\begin{lemma}[Local AD at one scale; WZ1, Lemma 4.11]
\label{lem:grain-ad-one-scale}
\lean{Kakeya.Assouad.WZ1LocalGrainOneScaleConclusion}
\uses{def:ad-set,def:extremal,def:plane-map}
For $\rho\in[\delta^{1-\eps},\delta^\eps]$, there is an
$\eps$-extremal refinement such that the projection of its shaded
union in each occupied $\sqrt\rho$-cube, onto the normal at that
cube, is a $(\rho,1-\sigma,\delta^{-\eps})_1$-ADset.
\end{lemma}
\begin{proof}
\uses{lem:grain-interval}
Apply the interval estimate at the required intermediate lengths.
At the smallest lengths use the trivial covering bound, and at the
largest lengths use the estimate for the entire cube. The parameter
choices absorb the losses from the finite list of lengths.
\end{proof}

\begin{lemma}[Local grains at every scale; WZ1, Lemma 4.12]
\label{lem:local-grains-all-scales}
\lean{Kakeya.Assouad.WZ1EveryScaleLocalGrainConclusion,Kakeya.Assouad.pure_wz2_every_scale_local_grain}
\uses{def:ad-set,def:extremal,def:plane-map}
Given $\eps>0$, a sufficiently extremal family with a sufficiently
slowly varying plane map has an $\eps$-extremal refinement for which
\[
        V(p)\cdot(E_{Y'}\cap B(p,\sqrt\rho))
\]
is a $(\rho,1-\sigma,\delta^{-\eps})_1$-ADset for every
$\rho\in[\delta,1]$ and every retained point $p$.
\end{lemma}
\begin{proof}
\uses{lem:grain-ad-one-scale}
Iterate the one-scale statement at a sufficiently fine finite list of
power scales, keeping the same restricted plane map. Upper AD bounds
survive each further refinement. Interpolate between the selected
scales using monotonicity of covering numbers; the scale spacing is
chosen so that this costs less than the remaining loss budget.
\end{proof}

\begin{remark}[A common configuration]
\label{rem:common-grains}
In Propositions~\ref{prop:lipschitz-grains} and
\ref{prop:c2-grains}, the global and local conclusions refer to one
family and one shading. Combining conclusions obtained from unrelated
refinements would not prove either proposition. The formalization
retains the same family, shading, density, volume bound, and
nonconcentration data through the joint selection. The local normal
also has to be transported as a covector under an affine map: if
$x'=Ax+b$, its new direction is proportional to $A^{-\mathsf T}V$.
Applying $A$ directly to the normal gives the wrong plane.
\end{remark}
